% fig_negseq_variability.tex — k₂ = |I₂|/|I₁| distribution: field vs synthetic model
\begin{tikzpicture}
\begin{axis}[
  thesis style,
  ybar interval,
  width  = 0.90\columnwidth,
  height = 4.8cm,
  xlabel = {Negative-sequence ratio $k_2 = |I_2|/|I_1|$},
  ylabel = {Count (LL and LLG faults, $n{=}50$)},
  xmin=-0.02, xmax=0.56,
  ymin=0, ymax=20,
  xtick={0,0.10,0.20,0.30,0.40,0.50},
  ytick={0,5,10,15,20},
  legend pos=north east,
  legend style={font=\scriptsize},
  grid=both,
  grid style={gray!20,line width=0.3pt},
]

%% Field measurements — continuous distribution
\addplot[fill=thesisblue!55,draw=thesisblue,opacity=0.8]
  coordinates {
    (0.00, 3)(0.05, 5)(0.10, 7)(0.15, 8)(0.20,11)(0.25,12)(0.30, 9)
    (0.35, 6)(0.40, 5)(0.45, 3)(0.50, 1)(0.55, 0)
  };
\addlegendentry{Field measurements ($k_2 \in [0.03,\,0.41]$)};

%% Synthetic model: two spikes at k₂=0 and k₂=0.5
%% Shown as vertical reference lines since the model is discrete
\addplot[fill=thesisgray!40,draw=thesisgray,opacity=0.6]
  coordinates {
    (0.00,18)(0.05,0)(0.45,0)(0.50,16)(0.55,0)
  };
\addlegendentry{Synthetic model ($k_2 \in \{0,\,0.5\}$)};

%% Annotations
\node[thesisblue,font=\tiny,align=center] at (axis cs:0.25,14)
  {Continuous field\\distribution};
\draw[thesisblue,->] (axis cs:0.25,13.5) -- (axis cs:0.25,12.2);

\node[thesisgray,font=\tiny,align=center] at (axis cs:0.43,17)
  {Model spike\\at $k_2{=}0.5$};

\end{axis}
\end{tikzpicture}
